% math-spine/expressive-manifolds.tex
% Shared-Horizon Preprint — Mathematical Spine
% Module 2: Expressive Manifolds

\section{Expressive Manifolds}
\label{sec:expressive-manifolds}

\subsection{Expressive State Space}

Let $\mathcal{E}$ denote the expressive state space associated with ANIMA and
Reversible Narrative Ecology (RNE). We model $\mathcal{E}$ as a smooth manifold
equipped with an expressive metric $g_{\mathcal{E}}$ capturing curvature,
affective gradients, and narrative coherence.

\[
(\mathcal{E}, g_{\mathcal{E}}) \quad \text{is a smooth expressive manifold.}
\]

The metric $g_{\mathcal{E}}$ is lineage-dependent but must satisfy:

\begin{enumerate}
    \item Smoothness: $g_{\mathcal{E}} \in C^\infty(T\mathcal{E} \times T\mathcal{E})$.
    \item Curvature boundedness: $|K(p)| < \infty$ for all $p \in \mathcal{E}$.
    \item Stability: expressive curvature cannot increase under valid transitions.
\end{enumerate}

\subsection{Expressive Curvature}

Expressive curvature $K$ measures the deviation of expressive trajectories from
geodesic behavior. For any point $p \in \mathcal{E}$:

\[
K(p) = \text{Curv}_{g_{\mathcal{E}}}(p).
\]

For ANIMA, curvature reflects phenomenological intensity.  
For RNE, curvature reflects narrative reversibility constraints.

\subsection{Expressive Transitions}

An expressive transition is a smooth map
\[
T : \mathcal{E} \rightarrow \mathcal{E}
\]
that preserves expressive curvature:

\[
K(T(p)) \le K(p) \quad \forall p \in \mathcal{E}.
\]

This condition ensures expressive stability and prevents collapse into operator
geometry (Section \ref{sec:operator-geometry}).

\subsection{Expressive Geodesics}

Geodesics on expressive manifolds represent minimal-change expressive
trajectories. A curve $\gamma : [0,1] \rightarrow \mathcal{E}$ is expressive
geodesic if:

\[
\nabla_{\dot{\gamma}} \dot{\gamma} = 0.
\]

Geodesics are used to compare expressive states across lineages without
activating narrative or operator surfaces.

\subsection{Expressive–Operator Separation}

To prevent collapse between expressive and operator geometries, we require:

\[
O \circ T \notin \text{Domain}(O)
\]
for any operator $O$ and expressive transition $T$.

This ensures that expressive manifolds remain distinct from operator manifolds
and preserves lineage sovereignty.

\subsection{Expressive Manifold Summary}

Expressive manifolds provide the geometric foundation for ANIMA and RNE. They
define curvature, stability, and geodesic behavior for expressive and narrative
systems. These manifolds must remain smooth, curvature-bounded, and separated
from operator geometry to maintain conceptual and mathematical integrity.
